\chapter{Pneumonia}

\section*{Definition}
Pneumonia is an acute infection of the lung parenchyma, diagnosed clinically by a constellation of symptoms (cough, fever, dyspnea, pleuritic pain) with confirmatory chest imaging demonstrating consolidation. The CURB-65 score (0--5) assesses severity and guides inpatient versus outpatient management.

\section*{What Happens in Your Body}
Pathogens reach the alveoli through microaspiration of oropharyngeal flora, inhalation of aerosolized droplets, or hematogenous spread. The host immune response triggers alveolar-capillary inflammation with neutrophil recruitment, leading to fibrinosuppurative consolidation, impaired gas exchange, and systemic inflammatory response.

\section*{Causes}
\begin{itemize}
  \item \textbf{Community-acquired:} \textit{Streptococcus pneumoniae} (most common), \textit{Mycoplasma pneumoniae}, \textit{Haemophilus influenzae}, \textit{Chlamydia pneumoniae}, \textit{Legionella pneumophila}, respiratory viruses (influenza, RSV, SARS-CoV-2)
  \item \textbf{Hospital-acquired:} \textit{Pseudomonas aeruginosa}, MRSA, \textit{Enterobacteriaceae}, \textit{Acinetobacter} species
  \item \textbf{Aspiration:} Anaerobic oral flora (\textit{Bacteroides}, \textit{Fusobacterium}), \textit{Peptostreptococcus}, often polymicrobial
  \item \textbf{Immunocompromised:} \textit{Pneumocystis jirovecii}, \textit{Aspergillus}, CMV, nontuberculous mycobacteria
\end{itemize}

\section*{When to See a Doctor}
See your doctor if:
\begin{itemize}
  \item Fever greater than 102\textdegree{}F (38.9\textdegree{}C) with chills and rigors
  \item Respiratory rate greater than 24 breaths/min, oxygen saturation less than 92\%, or use of accessory muscles
  \item Confusion, hypotension (systolic BP less than 90 mmHg), or signs of sepsis
  \item Immunocompromised state, age greater than 65, or multiple comorbidities
  \item Hemoptysis, chest pain, or inability to maintain oral hydration
\end{itemize}

\section*{Related Symptoms}
\begin{itemize}
  \item Cough (productive or dry), fever, dyspnea, pleuritic chest pain, sputum production, chills, fatigue, malaise
\end{itemize}
\textbf{Important:} Atypical pneumonia (\textit{Mycoplasma}, \textit{Chlamydia}, \textit{Legionella}) often presents with headache, myalgias, and dry cough with minimal chest findings, earning the historical name ``walking pneumonia.''

\section*{Self-Care / Home Management}
\begin{itemize}
  \item Complete prescribed antibiotic course fully; do not stop early even if symptoms improve
  \item Acetaminophen or ibuprofen for fever and pleuritic chest pain
  \item Increase fluid intake (water, warm tea, broth) to prevent dehydration and thin secretions
  \item Deep breathing and incentive spirometry to prevent atelectasis and aid secretion clearance
  \item Rest in semi-upright position; avoid supine lying to prevent secretion pooling
\end{itemize}

\section*{How Your Doctor Will Evaluate You}
\begin{itemize}
  \item Chest X-ray (PA and lateral) to confirm consolidation, characterize lobar vs. bronchopneumonia pattern
  \item Pulse oximetry and arterial blood gas for hypoxia assessment and severity stratification
  \item CBC with differential, serum chemistry, procalcitonin (elevated in bacterial pneumonia), blood cultures
  \item Sputum Gram stain and culture for hospitalized people; molecular panel (PCR) for atypical pathogens
  \item CT chest for complicated pneumonia (empyema, abscess, necrotizing pneumonia) or non-resolving cases
\end{itemize}

\section*{Treatment Options}
\begin{itemize}
  \item Antibiotics based on the most likely cause based on setting and CURB-65 score: macrolide (azithromycin) or doxycycline for mild CAP; beta-lactam plus macrolide for moderate severity
  \item Respiratory support: supplemental oxygen to maintain SpO2 greater than 92\%; noninvasive ventilation for hypercapnia
  \item Antiviral therapy (oseltamivir for influenza, remdesivir for COVID-19) when indicated
  \item Hospital admission for CURB-65 score of 2 or greater; ICU for septic shock or respiratory failure
  \item Drainage of parapneumonic effusion or empyema if present; chest tube for loculated collections
\end{itemize}

\section*{References}
\begin{thebibliography}{99}
\bibitem{pneumonia_symptoms:pneumonia1} Mandell LA, Wunderink RG, Anzueto A, et al. ``Infectious Diseases Society of America/American Thoracic Society consensus guidelines on the management of community-acquired pneumonia in adults.'' Clin Infect Dis. 2007;44(Suppl 2):S27-S72.
\bibitem{pneumonia_symptoms:pneumonia2} Musher DM, Thorner AR. ``Community-acquired pneumonia.'' N Engl J Med. 2014;371(17):1619-1628.
\bibitem{pneumonia_symptoms:pneumonia3} Jain S, Self WH, Wunderink RG, et al. ``Community-acquired pneumonia requiring hospitalization among U.S. adults.'' N Engl J Med. 2015;373(5):415-427.
\bibitem{pneumonia_symptoms:pneumonia4} File TM Jr, Marrie TJ. ``Clinical evaluation and diagnosis of pneumonia.'' UpToDate, 2023.
\bibitem{pneumonia_symptoms:pneumonia5} Wunderink RG, Waterer GW. ``Community-acquired pneumonia.'' In: Harrison\'s Principles of Internal Medicine. 21st ed. McGraw-Hill; 2022.
\bibitem{pneumonia_symptoms:pneumonia6} Metlay JP, Waterer GW, Long AC, et al. ``Diagnosis and treatment of adults with community-acquired pneumonia. An official clinical practice guideline of the ATS and IDSA.'' Am J Respir Crit Care Med. 2019;200(7):e45-e67.
\end{thebibliography}
